\chapter{Osteoporosis Screening --- DEXA Scan}

\section*{What Is This Test?}
This test measures specific substances in your blood, urine, or other body fluids to help diagnose, monitor, or screen for medical conditions. Your doctor will explain why this particular test was chosen for your situation.

\section*{Alternative Names}
DEXA scan, DXA, bone density test, dual-energy X-ray absorptiometry, bone densitometry, central DXA, peripheral DXA


\section*{Principle}
DXA uses two X-ray beams at different energies (typically 70 and 140 kVp). The differential attenuation of the two energies distinguishes bone from soft tissue. Bone mineral density (BMD) in g/cm2 is calculated. The T-score compares the patient's BMD to a young adult reference population.

\section*{History}
Bone densitometry evolved from plain X-ray techniques in the 1960s. Dual-energy X-ray absorptiometry (DXA) was developed in the 1980s and became the gold standard for osteoporosis diagnosis after the WHO defined T-score criteria in 1994.

\section*{Why This Test Is Ordered}
Your bones get thinner as you age, especially after menopause. A DEXA scan measures bone density to find out if you are at risk for fractures before one actually happens. It answers the question: ``Are your bones strong enough to survive a fall?''

\section*{Is This a Primary or Secondary Test}
Primary screening test for osteoporosis.

\section*{How This Test Is Performed}
You lie on a padded table while a machine scans your hip and spine with a low-dose X-ray. The scan takes about 15 minutes. There is no needle, no injection, and no pain. You stay fully clothed.

\section*{What Preparation Is Needed}
Some facilities ask patients to avoid calcium supplements for a period before DXA; follow local appointment instructions rather than stopping a prescribed supplement on your own. Wear clothing without metal fasteners when requested, and tell the imaging team if you are pregnant or may be pregnant.

\section*{Contraindications and Precautions}
DEXA scanning uses very low-dose X-rays. It is generally not performed during pregnancy. Inform the team if you are pregnant or may be pregnant. There are no other specific contraindications.

\section*{Risks}
DEXA scanning uses very low-dose ionizing radiation, much less than a standard X-ray. The radiation exposure is minimal and the benefit of assessing fracture risk outweighs the risk. Inform the team if you are pregnant. There are no other significant risks.

\section*{What Happens After the Test}
You can resume normal activities immediately. Results usually take a few days. Your doctor will compare your bone density to two numbers: the peak bone density of a healthy 30-year-old (your T-score) and the average for someone your age and sex (your Z-score).

\section*{Understanding Your Results}

\section*{Normal Results}
Normal: T-score $\ge$ -1.0. Osteopenia: T-score -1.0 to -2.5. Osteoporosis: T-score $\le$ -2.5.

\section*{Follow-up Tests}
If the DEXA scan shows low bone density, your doctor may order blood tests (calcium, vitamin D, thyroid function, bone markers), a vertebral fracture assessment, or repeat DEXA scanning in 1--2 years to monitor changes.
\clearpage
